\documentclass{article}
\usepackage{amsmath}
\usepackage{amssymb}
\usepackage{amsthm}
\usepackage{enumitem}
\usepackage{tikz-cd}
\usepackage[a4paper, margin=1in]{geometry}

\newtheorem{theorem}{Teorema}
\newtheorem{corollary}{Corollario}[theorem]

\begin{document}
\setlength{\parindent}{0pt}

\section{Teoremi ausiliari}
    \subsection{Generatori, basi e dimensione di uno spazio vettoriale}
    Un sottoinsieme $G$ di uno spazio vettoriale $V$ che sussiste su un campo $F$ è definito generatore di $V$ se e solo se $\textnormal{span}(G) = V$, dove $\textnormal{span}(G) = \{\mathbf{v} \in V \; | \; \displaystyle\sum_{i = 1}^{|G|} a_i \mathbf{g_i} = \mathbf{v}, \; \textnormal{con} \; a_1, ..., a_{|G|} \in F\}$.

    \vspace{10pt}
    Un sottoinsieme $B = \{\mathbf{b_1},...,\mathbf{b_n}\}$ di uno spazio vettoriale $V$ è definito base se e solo se:
    \begin{enumerate}[label=(\arabic*)]
        \item $ \forall \mathbf{v} \in V, \; \exists \, c_1, ..., c_n \in F \; | \; \displaystyle\sum_{i = 1}^{n} c_i \mathbf{b_i} = \mathbf{v}, \quad \mathbf{b_1}, ..., \mathbf{b_n} \in B$ (in altre parole $B$ è un generatore di $V$) 
        \item Per ogni sottoinsieme $\{\mathbf{b_1}, ..., \mathbf{b_k}\} \subset B$ , $c_1 \mathbf{b_1}, ..., c_k \mathbf{b_k} = \mathbf{0} \iff c_1 = c_2 = ... = c_k = 0 $  
    \end{enumerate}

    \vspace{5pt}
    La dimensione $n$ di uno spazio vettoriale è data dalla cardinalità di una qualsiasi base $B$ di $V$.
    Inoltre, dalla definizione di base segue che, dato un vettore $\mathbf{v}$ e una base $B$ di $V$, vi è una e una sola scomposizione in coordinate secondo $B$ di $\mathbf{v}$.

    \vspace{10pt}
    \begin{theorem}[Teorema dell'unicità della scomposizione]
        Sia $V$ uno spazio vettoriale su $F$ e sia $\{\mathbf{b_1},...,\mathbf{b_n}\}$ una base di $V$. Allora per ogni vettore $\mathbf{v} \in V$ esiste una e una sola combinazione
        $c_1,...,c_n \in F$ tale per cui $c_1\mathbf{b_1} + ... + c_n\mathbf{b_n} = \mathbf{v}$
    \end{theorem}

    \begin{proof}[Dimostrazione]
        Supponiamo per assurdo che vi siano due combinazioni $a_1,...,a_n \in F$ e $b_1,...,b_n \in F$ tali per cui: 
        \begin{enumerate}[label=(\arabic*)]
            \item $a_i \neq b_i \; \forall i \in \{1,...,n\}$
            \item $a_1\mathbf{b_1} + ... + a_n\mathbf{b_n} = \mathbf{v}$ 
            \item $b_1\mathbf{b_1} + ... + b_n\mathbf{b_n} = \mathbf{v}$ 
        \end{enumerate} 

        Allora, segue che:
        \[
            (a_1\mathbf{b_1} - b_1\mathbf{b_1}) + ... + (a_n\mathbf{b_n} - b_n\mathbf{b_n}) = \mathbf{v} - \mathbf{v} = \mathbf{0} \implies (a_1 - b_1) \mathbf{b_1} + ... + (a_n - b_n) \mathbf{b_n} = \mathbf{0} 
        \]

        Per il punto (2) della definizione di base, abbiamo dunque:
        \[
            (a_i - b_i) = 0 \implies a_i = b_i \quad \forall i \in \{1,...,n\}
        \]
        chiaramente contraddicendo l'ipotesi (1).
    \end{proof}

    \begin{theorem}[Teorema dell'immagine di base]
        Sia $f : V \rightarrow W$ un'applicazione lineare e sia $B = \{\mathbf{b_1}, ..., \mathbf{b_n}\}$ una base di $V$, allora $f(B) = \{f(\mathbf{b_1}), ..., f(\mathbf{b_n})\}$ è un generatore di $\textnormal{Im}(f)$.
    \end{theorem}

    \begin{proof}[Dimostrazione]
        Si prenda un $\mathbf{w} \in \textnormal{Im}(f)$, allora $\exists \mathbf{v} \in V$ tale cui $f(\mathbf{v}) = \mathbf{w}$. Dato che $B$ è una base di $V$, si ha:
        \[
            \mathbf{w} = f(\mathbf{v}) = f(\sum_{i} a_i \mathbf{b_i}) = \sum_{i} f(a_i \mathbf{b_i}) = \sum_{i} a_i f(\mathbf{b_i}) \implies \textnormal{Im}(f) \subset \textnormal{Span}(\{f(\mathbf{b_1}), ..., f(\mathbf{b_n})\})
        \]

        Sia $\mathbf{k} \in \textnormal{Span}(\{f(\mathbf{b_1}), ..., f(\mathbf{b_n})\})$, allora:
        \[
            \mathbf{k} = \sum_{i} c_i f(\mathbf{b_i}) = f(\sum_{i} c_i \mathbf{b_i}) \implies \exists \mathbf{u} = \sum_{i} c_i \mathbf{b_i} \in V \; | \; f(\mathbf{u}) = \mathbf{k} \implies \textnormal{Span}(\{f(\mathbf{b_1}), ..., f(\mathbf{b_n})\}) \subset \textnormal{Im}(f)
        \]

        Quindi $\textnormal{Im}(f) = \textnormal{Span}(\{f(\mathbf{b_1}), ..., f(\mathbf{b_n})\})$. 
    \end{proof}

\end{document}